% ID: SEL-D0058-OA-E02
% ESTADO: validada. S058.
\textbf{a)} El denominador se anula sólo en $x=-1/2$, así que el dominio es $\mathbb R\setminus\{-1/2\}$. El único cero del numerador es $x=0$, donde el denominador vale $1$. Por tanto, ambos cortes con los ejes coinciden en $(0,0)$. Como
\[
f(x)=2-\frac{2}{2x+1},
\]
hay asíntota vertical $x=-1/2$ y horizontal $y=2$. El término fraccionario diverge a ambos lados del polo y tiende a cero en ambos infinitos.

\textbf{b)} $g'(x)=3(x+1)^2$, positiva salvo en $x=-1$. Como $g(x)=(x+1)^3-1$, es estrictamente creciente en todo $\mathbb R$. En $-1$ se anula la derivada sin cambiar de signo y no hay extremos relativos. La segunda derivada es $g''(x)=6(x+1)$: la gráfica es cóncava hacia abajo en $(-\infty,-1)$ y hacia arriba en $(-1,\infty)$. Cambia de curvatura en $\boxed{(-1,-1)}$, con tangente horizontal.
